\documentclass[crop,tikz,11pt]{standalone}
\usepackage{geometricfigures}
\usepackage{amsmath,amssymb}
\usepackage{pgfplots}
\pgfplotsset{compat=1.18}
\usetikzlibrary{arrows.meta,decorations.markings,patterns,intersections}
\usepgfplotslibrary{fillbetween}

\definecolor{lib}{HTML}{3B75AF}
\definecolor{rot}{HTML}{C1701E}
\definecolor{sep}{HTML}{B03030}
\definecolor{net}{HTML}{4E9A4E}

\begin{document}
\begin{tikzpicture}[baseline]
\begin{axis}[width=7.4cm,height=5.2cm,xmin=0,xmax=6.2832,ymin=-2.6,ymax=2.6,
  xtick={0,3.1416,6.2832},xticklabels={$0$,$\pi$,$2\pi$},
  xlabel={$\theta$},ylabel={$p_\theta$},axis lines=left,title style={at={(0.5,0)},anchor=north,yshift=-3.5em},   % panel label below the panel
  title={\small librating: area \emph{enclosed}},
  every axis plot/.append style={smooth,samples=200}]
  % H = 0.5 turns at theta = arccos(0.5) = pi/3 and 2pi - pi/3, and the curve is
  % drawn only between them. Running the domain to 0:2pi instead would leave the
  % clamped sqrt sitting at zero outside the loop and draw a spurious horizontal
  % line along p_theta = 0 -- see the note in pendulum-phase-portrait.
  \addplot[lib,thick,name path=U,domain=1.0472:5.2360]
    { sqrt(max(0,2*(0.5-cos(deg(x)))))};
  \addplot[lib,thick,name path=L,domain=1.0472:5.2360]
    {-sqrt(max(0,2*(0.5-cos(deg(x)))))};
  \addplot[lib!25!bg]fill between[of=U and L];
  \addplot[sep,dashed,domain=0:6.2832]{ 2*abs(sin(deg(x/2)))};
  \addplot[sep,dashed,domain=0:6.2832]{-2*abs(sin(deg(x/2)))};
\end{axis}
\end{tikzpicture}
\end{document}
